% !TeX root = main.tex
% Reusable manuscript scaffold; see README.md (Japanese).
% Search for TODO and replace_me before submission.
\RequirePackage{iftex}
\ifPDFTeX
  % pdfLaTeX selects its graphics driver automatically.
\else
  % This branch supports the original pLaTeX + dvipdfmx workflow.
  % LuaLaTeX and XeLaTeX require separate engine/driver configuration.
  \PassOptionsToPackage{dvipdfmx}{graphicx}
  \PassOptionsToPackage{dvipdfmx}{hyperref}
  \PassOptionsToPackage{dvipdfmx}{xcolor}
  \PassOptionsToPackage{dvipdfmx}{color}
  \PassOptionsToPackage{dvipdfmx}{geometry}
\fi
\documentclass[a4paper,fleqn]{cas-sc}
\usepackage[authoryear,longnamesfirst]{natbib}
\usepackage{amsmath}
\usepackage{mathtools}
\usepackage{bm}
\usepackage{subcaption}
\usepackage{algorithm}
\usepackage{algpseudocode}
\usepackage{amssymb,pifont,booktabs,multirow}
\usepackage{rotating}
\usepackage{placeins}
\usepackage{threeparttable}

\newcommand{\cmark}{\checkmark}
\newcommand{\xmark}{\ding{55}}
\newcommand{\mmark}{\ensuremath{\triangle}}
\newcommand{\pmark}{\mmark}

% Compatibility and caption-width fixes retained from the source manuscript.
\AddToHook{begindocument/before}{\ifdefined\bibsep\else\newlength\bibsep\fi}
\ExplSyntaxOn
\exp_args:NNno \exp_args:Nno \use:n { \cs_gset:Npn \__make_fig_caption:nn #1#2 }
  {
    \exp_after:wN \use_ii_i:nn \exp_after:wN
      { \__make_fig_caption:nn {#1} {#2} }
      { \dim_set:Nn \l_fig_width_dim \linewidth }
  }
\exp_args:NNno \exp_args:Nno \use:n { \cs_gset:Npn \__make_tbl_caption:nn #1#2 }
  {
    \exp_after:wN \use_ii_i:nn \exp_after:wN
      { \__make_tbl_caption:nn {#1} {#2} }
      { \dim_set:Nn \l_tbl_width_dim \linewidth }
  }
% Avoid an empty ORCID footnote when no verified identifiers are supplied.
\cs_new_eq:NN \template_printorcid: \printorcid
\RenewDocumentCommand \printorcid { }
  { \seq_if_empty:NF \g_stm_orcid_seq { \template_printorcid: } }
\ExplSyntaxOff

% Template-only helpers. Delete these after replacing their uses.
\newcommand{\placeholder}[1]{\textcolor{gray}{\textit{[TODO: #1]}}}
\newcommand{\figureplaceholder}[2]{%
  \begingroup
  \setlength{\fboxsep}{10pt}%
  \fbox{\parbox[c][#1][c]{\dimexpr\linewidth-2\fboxsep-2\fboxrule\relax}{%
    \centering\sffamily\color{gray}#2}}%
  \endgroup
}
\graphicspath{{figures/}{./}}

% Set false to omit these optional parts. Check the destination journal.
\newif\ifshowgraphicalabstract
\showgraphicalabstracttrue
\newif\ifshowhighlights
\showhighlightstrue
\newif\ifshowformattingexamples
\showformattingexamplestrue

\begin{document}
\let\WriteBookmarks\relax
\def\floatpagepagefraction{1}
\def\textpagefraction{.001}

% ===== 1. Title, authors, affiliations =====
\shorttitle{Short title of the new manuscript} % TODO
\shortauthors{F. Author and S. Author} % TODO
\title[mode = title]{Title of the new manuscript} % TODO

\author[1]{First Author} % TODO: add [orcid=...] only after verification.
\ead{first.author@example.com} % TODO
% \credit{...} % TODO: enter agreed author contributions, if applicable.

\affiliation[1]{
  organization={Department, University or Institution}, % TODO
  addressline={Street address}, % TODO
  city={City}, % TODO
  postcode={Postal code}, % TODO
  state={State or region}, % TODO
  country={Country} % TODO
}

\author[1]{Second Author} % TODO
\cormark[1] % TODO: assign to the actual corresponding author.
\cortext[1]{Corresponding author}
\ead{second.author@example.com} % TODO
% \credit{...}
% For another affiliation, use \author[2]{...} and \affiliation[2]{...}.

% ===== 2. Abstract and optional front matter =====
% Explicit optional argument avoids losing the first token in CAS's verbatim abstract.
\begin{abstract}[]
\placeholder{State the research problem and objective. Describe the method and
the evidence used to evaluate it. Report only supported findings, with the
relevant scope and uncertainty. End with the contribution and its limitations.
Replace this entire placeholder with the abstract of the new manuscript.}
\end{abstract}

\ifshowgraphicalabstract
\begin{graphicalabstract}
\centering
% Replace the box with, for example:
% \includegraphics[width=\linewidth]{graphical-abstract.pdf}
\figureplaceholder{48mm}{Graphical abstract\\[4pt]TODO: replace with your own figure}
\end{graphicalabstract}
\fi

\ifshowhighlights
\begin{highlights}
  \item \placeholder{Main contribution, within the journal's length limit.}
  \item \placeholder{Distinctive feature of the method.}
  \item \placeholder{Finding supported by the new study's evidence.}
\end{highlights}
\fi

\begin{keywords}
Keyword one \sep Keyword two \sep Keyword three % TODO
\end{keywords}
\maketitle

% ===== 3. Main text: retain or adapt the section functions =====
\section{Introduction}
\label{sec:introduction}
\placeholder{Explain the problem, why it matters, and the limitation in existing
knowledge. State the research question and the contribution supported by this
paper. Describe the organization if it helps the reader.}

\section{Literature review}
\label{sec:literature}
\placeholder{Organize relevant studies by the scientific distinctions needed to
understand the research gap. Verify each citation and explain the paper's position.}

% Formatting example only. The following citation is deliberately fictitious.
Author--year citation examples: \citet{replace_me_reference} and
\citep{replace_me_reference}. Replace both the citation and its BibTeX entry.

\section{Methodology}
\label{sec:methodology}
\subsection{Problem statement}
\label{subsec:problem}
\placeholder{Define the system, notation, decisions, assumptions, and research
scope. Distinguish assumptions from consequences derived from the model.}

\subsection{Model or method}
\label{subsec:method}
\placeholder{Describe the model, estimation, optimization, or analytical method
in sufficient detail to reproduce it. Include only justified mathematical claims.}

\section{Numerical experiments}
\label{sec:experiments}
\subsection{Data and experimental settings}
\label{subsec:settings}
\placeholder{Document data sources, scenarios, baselines, evaluation metrics,
parameter settings, uncertainty, and reproducibility information. Adapt this
section if the paper uses empirical analysis or theoretical examples instead.}

\subsection{Results and interpretation}
\label{subsec:results}
\placeholder{Present the evidence for each claim. Separate observed results
from explanations and generalizations. Discuss uncertainty and limitations.}

% Optional: add a separate Discussion section when needed.
% \section{Discussion}
% \label{sec:discussion}

\section{Conclusion}
\label{sec:conclusion}
\placeholder{Answer the research question using the supported findings.
State the contribution, limits of applicability, and justified future work.}

% ===== 4. End matter =====
% Uncomment after entering verified \credit{...} statements for the authors.
% \printcredits

\section*{Acknowledgments}
\placeholder{Enter only funding identifiers and acknowledgments that apply to
the new study. Remove this section if it is not applicable.}

% Add required declarations only after checking the destination journal and
% actual circumstances. These statements must not be copied from another study.
% \section*{Declaration of competing interest}
% \section*{Data availability}
% \section*{Declaration of generative AI and AI-assisted technologies
% in the manuscript preparation process}

% ===== 5. Appendices =====
\appendix
\section{Supporting material}
\label{app:support}
\placeholder{Place derivations, additional results, or reproducibility details
here when they are needed. Remove unused appendices.}

\ifshowformattingexamples
\clearpage
\section{Formatting examples (delete before submission)}
\label{app:formatting}
The following objects demonstrate formatting only. They do not state a model,
an algorithm, or a numerical finding for the new study.

\subsection{Equations and mathematical notation}
\begin{equation}
  y = f(\bm{x};\bm{\theta}).
  \label{eq:formatting-example}
\end{equation}
Equation~\eqref{eq:formatting-example} demonstrates left-aligned mathematics,
bold symbols, and an equation reference. Replace the expression and define
every symbol used in the actual study.

\subsection{Figures and subfigures}
\begin{figure}[pos=htbp]
  \centering
  \begin{subfigure}{0.47\linewidth}
    \centering
    \figureplaceholder{25mm}{Panel A: TODO}
    % \includegraphics[width=\linewidth]{panel-a.pdf}
    \caption{First panel.}
    \label{fig:panel-a}
  \end{subfigure}\hfill
  \begin{subfigure}{0.47\linewidth}
    \centering
    \figureplaceholder{25mm}{Panel B: TODO}
    % \includegraphics[width=\linewidth]{panel-b.pdf}
    \caption{Second panel.}
    \label{fig:panel-b}
  \end{subfigure}
  \caption{Formatting example: two panels with a shared caption. Replace the
  boxes with the study's own figures, and explain the quantities and units.}
  \label{fig:formatting-example}
\end{figure}
Figure~\ref{fig:formatting-example} illustrates the caption and reference style.
Put new figure files in \texttt{figures/}.

\subsection{Tables and notes}
\begin{table}[pos=htbp]
  \centering
  \caption{Formatting example only; no numerical results are supplied.}
  \label{tab:formatting-example}
  \begin{threeparttable}
    \begin{tabular}{lcc}
      \toprule
      Method & Metric A (unit) & Metric B (unit) \\
      \midrule
      Baseline & --- & --- \\
      Proposed method & --- & --- \\
      \bottomrule
    \end{tabular}
    \begin{tablenotes}[flushleft]\footnotesize
      \item Note: Dashes are placeholders, not zero measurements.
    \end{tablenotes}
  \end{threeparttable}
\end{table}
Table~\ref{tab:formatting-example} uses the original table packages and style.

\FloatBarrier
\subsection{Algorithms}
\begin{algorithm}[htbp]
  \caption{Formatting example only: replace with the actual procedure}
  \label{alg:formatting-example}
  \begin{algorithmic}[1]
    \Require \textit{TODO: specify the inputs}
    \Ensure \textit{TODO: specify the outputs}
    \State \textit{TODO: describe initialization}
    \State \textit{TODO: describe each step and any stopping rule}
    \State \Return \textit{TODO: specify the returned quantity}
  \end{algorithmic}
\end{algorithm}
Algorithm~\ref{alg:formatting-example} demonstrates numbered pseudocode.
\FloatBarrier
\fi

\bibliographystyle{cas-model2-names}
\bibliography{references}
\end{document}
